% Main manuscript for Quantum Science and Technology / IOP submission.
% This draft uses standard LaTeX so it compiles without proprietary class files.
% IOP states that any common TeX variant is acceptable; if desired, replace
% the next line with \documentclass[anonymous]{iopjournal} or \documentclass[12pt]{iopart}.
\documentclass[12pt]{article}

\usepackage[a4paper,margin=1in]{geometry}
\usepackage{amsmath,amssymb,amsfonts,amsthm,mathtools}
\usepackage{bm}
\usepackage{graphicx}
\usepackage{booktabs}
\usepackage{algorithm}
\usepackage{algorithmic}
\usepackage{float}
\usepackage{enumitem}
\usepackage{microtype}
\usepackage{xcolor}
\usepackage[colorlinks=true,linkcolor=blue,citecolor=blue,urlcolor=blue]{hyperref}

\graphicspath{{figures/}}

\newtheorem{theorem}{Theorem}
\newtheorem{lemma}[theorem]{Lemma}
\newtheorem{proposition}[theorem]{Proposition}
\newtheorem{corollary}[theorem]{Corollary}
\newtheorem{definition}{Definition}
\newtheorem{assumption}{Assumption}
\newtheorem{remark}{Remark}

\DeclareMathOperator{\Tr}{Tr}
\DeclareMathOperator{\diag}{diag}
\DeclareMathOperator{\rank}{rank}
\DeclareMathOperator{\softplus}{softplus}
\newcommand{\R}{\mathbb{R}}
\newcommand{\C}{\mathbb{C}}
\newcommand{\calH}{\mathcal{H}}
\newcommand{\calJ}{\mathcal{J}}
\newcommand{\calD}{\mathcal{D}}
\newcommand{\calL}{\mathcal{L}}
\newcommand{\calM}{\mathcal{M}}
\newcommand{\calO}{\mathcal{O}}
\newcommand{\calR}{\mathcal{R}}
\newcommand{\calT}{\mathcal{T}}
\newcommand{\calE}{\mathcal{E}}
\newcommand{\id}{\mathrm{id}}
\newcommand{\dd}{\mathrm{d}}
\newcommand{\ii}{\mathrm{i}}
\newcommand{\eps}{\varepsilon}
\newcommand{\norm}[1]{\left\lVert #1 \right\rVert}
\newcommand{\abs}[1]{\left\lvert #1 \right\rvert}
\newcommand{\ip}[2]{\left\langle #1,#2 \right\rangle}
\newcommand{\ket}[1]{\left\lvert #1 \right\rangle}
\newcommand{\bra}[1]{\left\langle #1 \right\rvert}
\newcommand{\proj}[1]{\ket{#1}\!\bra{#1}}
\newcommand{\hilb}{\mathcal{H}}
\newcommand{\rhohat}{\widehat{\rho}}
\newcommand{\rhophi}{\rho_{\phi}}
\newcommand{\Thetahat}{\widehat{\Theta}}

\title{CPTP-Compiler-PINNs: Structure-Preserving Residual Learning and Matrix-Free Compilation for Open Quantum Dynamics}
\author{Your Name\\\small Independent Researcher / replace with affiliation for submission}
\date{\today}

\begin{document}
\maketitle

\begin{abstract}
Learning open quantum dynamics from sparse, noisy observations is a central task in quantum device characterization, noise modeling, and control. Standard neural surrogates for density matrices, however, can leave the physical state space during training, and dense Liouvillian implementations can materialize $d^2\times d^2$ superoperators even when the physical model is specified by only a few local Hamiltonian and dissipative terms. We introduce \emph{CPTP-Compiler-PINNs}, a structure-preserving residual-learning framework for Markovian open quantum systems. The method combines: (i) a Cholesky density network that maps every time point to a positive semidefinite, unit-trace density matrix; (ii) a softplus-positive Lindblad-rate parameterization, or a positive Kossakowski parameterization, that preserves the Gorini-Kossakowski-Sudarshan-Lindblad generator form; and (iii) a matrix-free compiler that lowers a symbolic open-system specification to either dense or structured residual kernels. We prove exact physicality of the density map, universal approximation of continuous density trajectories including boundary-rank states, a complete-positivity and trace-preservation guarantee for the learned generator, a trace-norm residual certificate for trajectory error, a local identifiability theorem based on the sensitivity Gramian, and a dense-versus-structured compiler complexity theorem. Authentic CPU experiments generated by the accompanying scripts validate the claims on a qubit inverse problem, a residual-certificate stress test, qutrit leakage identification, and dense-versus-structured residual benchmarking. The reported simulations are intentionally reproducible and modest in scale; they are not presented as hardware data or as state-of-the-art performance claims. The contribution is a mathematically constrained, accelerator-aware formulation that can be scaled with batched GPU kernels and integrated into quantum software stacks.
\end{abstract}

\noindent\textbf{Keywords:} open quantum systems; Lindblad master equation; physics-informed neural networks; completely positive trace-preserving maps; quantum system identification; matrix-free compilers; CUDA-Q; scientific machine learning

\tableofcontents

% ============================================================
\section{Introduction}
\label{sec:introduction}

Open quantum systems are the operational language of modern quantum hardware. A superconducting, trapped-ion, photonic, spin, or neutral-atom processor is never described only by its ideal Hamiltonian; it is also described by relaxation, dephasing, leakage, crosstalk, drift, and control-induced dissipation. The mathematical object that must remain physical throughout this description is the density matrix. For a finite Hilbert space $\hilb\cong\C^d$, the state space is
\begin{equation}
    \calD(\hilb)=\{\rho\in\C^{d\times d}:\rho=\rho^\dagger,\ \rho\succeq0,\ \Tr\rho=1\}.
\end{equation}
When the dynamics are Markovian and time-local, the most general generator of a completely positive trace-preserving (CPTP) semigroup has the Gorini-Kossakowski-Sudarshan-Lindblad (GKSL) form~\cite{lindblad1976generators,gorini1976completely,breuer2002theory}. This form is the backbone of quantum noise modeling, quantum control, and calibration workflows~\cite{wiseman2009quantum,krantz2019quantum,leibfried2003quantum,preskill2018quantum}.

The inverse problem is increasingly important: infer a physically admissible dynamical model from incomplete observations. Classical process tomography reconstructs a channel but scales poorly with system size~\cite{chuang1997prescription,poyatos1997complete,knee2018quantum}. Hamiltonian learning and Lindblad learning exploit dynamical structure but still face identifiability and optimization challenges~\cite{boulant2003robust,wang2017experimental}. Neural methods and physics-informed neural networks (PINNs) offer flexible function approximation with equation residuals in the training objective~\cite{raissi2019physics,karniadakis2021physics,chen2018neural}. Yet a direct neural map $t\mapsto \rho_\theta(t)$ usually has no reason to remain Hermitian, positive semidefinite, or unit trace. A loss penalty for negative eigenvalues is not equivalent to physicality during training; it merely asks the optimizer to repair an invalid output after the fact.

This paper takes the opposite design principle: \emph{make invalid states unrepresentable}. The proposed density network is
\begin{equation}
    \rhophi(t)=\frac{A_\phi(t)A_\phi(t)^\dagger}{\Tr[A_\phi(t)A_\phi(t)^\dagger]},
    \label{eq:intro_cholesky}
\end{equation}
where $A_\phi(t)$ is a neural lower-triangular complex matrix with nonzero diagonal. This Cholesky-type map is differentiable, architecture-level physical, and compatible with either automatic differentiation or stable finite-difference residuals. Dissipation rates are parameterized by softplus transforms, and full Kossakowski matrices can be represented as $BB^\dagger$. The learning objective combines data likelihood, initial-state matching, and a Lindblad residual evaluated by a compiler that can avoid explicit $d^2\times d^2$ Liouvillians.

The word ``compiler'' is used deliberately. A quantum model is specified in an intermediate representation (IR): Hilbert dimension, Hamiltonian terms, jump operators, rates, observables, batching rules, and differentiation mode. The compiler lowers this IR to either a dense Liouvillian or a structured residual kernel. For small $d$, dense kernels can be convenient. For larger $d$, the structured route stores $H$ and $m$ Lindblad operators, not a $d^2\times d^2$ matrix. This aligns with the design philosophy of contemporary quantum and accelerator software: represent intent at a high level, then lower to hardware-conscious kernels. In that sense, the framework is naturally compatible with PyTorch~\cite{paszke2019pytorch}, SciPy~\cite{virtanen2020scipy}, Qiskit~\cite{qiskit2024}, CUDA-Q~\cite{cudaq2023}, and future custom GPU kernels.

\subsection{Contributions}
\label{sec:contributions}

The draft originally supplied with this work contained a promising structure but made claims that depended on missing figure and table files. The present manuscript replaces those claims with reproducible results generated by the included scripts. The core contributions are:
\begin{enumerate}[leftmargin=1.5em]
    \item \textbf{A hard-constrained density network.} The Cholesky density map enforces Hermiticity, positive semidefiniteness, and unit trace exactly at every time point, every training iteration, and every evaluation point.
    \item \textbf{A CPTP generator parameterization.} Nonnegative Lindblad rates are represented by softplus parameters; full Kossakowski matrices are represented as positive semidefinite Gram matrices. The learned generator therefore remains in GKSL form.
    \item \textbf{A matrix-free open-system compiler.} A symbolic open-system IR is lowered to either a dense superoperator or a structured residual kernel. The theorem in Section~\ref{sec:compiler_theory} quantifies memory and arithmetic cost.
    \item \textbf{Six mathematical guarantees.} We prove physicality, universal trajectory approximation, CPTP semigroup preservation, a residual-to-error certificate in trace norm, a local identifiability bound, and compiler complexity.
    \item \textbf{Authentic numerical validation.} The figures and tables are generated by scripts in the accompanying folder. They include a qubit finite-difference PINN inverse problem, a residual certificate stress test, qutrit leakage identification, and dense-versus-structured residual benchmarks.
\end{enumerate}

The results are intentionally stated conservatively. They show the value of the structural constraints and the compiler path on controlled simulations. They do not claim quantum-hardware validation, universal superiority over all baselines, or production-scale GPU performance. Those are future experimental targets, not present facts.


% ============================================================
\section{Related work and positioning}
\label{sec:related_work}

\subsection{Open-system identification and tomography}

The traditional route to learning quantum dynamics is tomography. Quantum process tomography estimates a channel from tomographically complete preparations and measurements~\cite{chuang1997prescription,poyatos1997complete}. It is model agnostic, but the number of settings grows rapidly with Hilbert-space dimension. Completely positive projection methods can repair an estimated map, but a projection step is not the same as learning inside the physical set throughout optimization~\cite{knee2018quantum}. Lindblad identification methods exploit the generator structure and can reduce the number of parameters when the dissipative mechanisms are known or partially known~\cite{boulant2003robust}. Hamiltonian learning similarly exploits a structured dynamical model and has been demonstrated experimentally~\cite{wang2017experimental}. The present method follows the structured-identification philosophy but replaces unconstrained trajectory estimation by a hard density-matrix parameterization.

\subsection{Physics-informed machine learning}

PINNs were introduced as a general framework for solving forward and inverse problems by penalizing differential-equation residuals at collocation points~\cite{raissi2019physics}. Their success depends on optimization, sampling, scaling of loss terms, and causal or curriculum training strategies~\cite{karniadakis2021physics,wang2022respecting}. The neural tangent kernel literature helps explain some optimization behavior in overparameterized networks~\cite{jacot2018neural}, while neural ODEs provide a complementary viewpoint in which the differential equation itself is a learnable continuous-depth model~\cite{chen2018neural}. Operator-learning models such as DeepONet and Fourier neural operators target families of solution maps rather than a single trajectory~\cite{lu2021learning,li2021fourier}. CPTP-Compiler-PINNs are closest to inverse PINNs: the network represents a trajectory, and generator parameters are inferred by matching both data and physics residuals.

\subsection{Quantum neural representations}

Neural quantum states have become a powerful representation for many-body wavefunctions and density operators~\cite{carleo2017solving,torlai2018neural}. Variational neural ansatzes have also been applied to open-system steady states~\cite{yoshioka2020variational}. Quantum machine learning more broadly studies quantum feature spaces, hybrid algorithms, and variational circuits~\cite{mcclean2016theory,schuld2019quantum}. These works are complementary. The present manuscript does not propose a quantum circuit ansatz or a many-body wavefunction ansatz. It proposes a constrained classical surrogate for density-matrix trajectories and a compiler path for evaluating Lindblad residuals efficiently.

\subsection{Dissipation as a resource and non-Markovian extensions}

Dissipation is not only a nuisance; it can be engineered for state preparation and quantum computation~\cite{verstraete2009quantum}. Open quantum walks and non-Markovian models show that richer dynamical structures may be necessary outside the GKSL semigroup setting~\cite{denisov2019open,wolf2008assessing,banchi2018modelling,luchnikov2020machine}. The method here deliberately starts with the Markovian case because it is the right abstraction for many calibration and noise-modeling tasks. The same density parameterization could be used inside a non-Markovian embedding, but the generator theorem would need to be replaced by a theorem for the enlarged model.

% ============================================================
\section{Background}
\label{sec:background}

\subsection{GKSL dynamics}
\label{sec:gksl}

Let $\rho(t)\in\calD(\hilb)$ be a density matrix. A time-independent Markovian master equation has the form
\begin{equation}
    \frac{\dd\rho}{\dd t}=\calL_\Theta[\rho]
    =-\ii[H_\Theta,\rho]+
    \sum_{k=1}^{m}\gamma_k\left(L_k\rho L_k^\dagger-\frac12\{L_k^\dagger L_k,\rho\}\right),
    \label{eq:gksl_rates}
\end{equation}
where $H_\Theta=H_\Theta^\dagger$, $\gamma_k\ge0$, and the $L_k$ are jump operators. The anticommutator is $\{A,B\}=AB+BA$. A time-dependent generator $\calL_t$ can also be used provided it remains in GKSL form at each time. The original GKSL theorem states that this structure characterizes generators of quantum dynamical semigroups in finite dimensions~\cite{lindblad1976generators,gorini1976completely}. Standard treatments of open quantum systems and quantum measurement build on this fact~\cite{breuer2002theory,wiseman2009quantum}.

For vectorization, we use column-major convention $\operatorname{vec}(A\rho B)=(B^T\otimes A)\operatorname{vec}(\rho)$. The dense Liouvillian associated with~\eqref{eq:gksl_rates} is
\begin{align}
    \mathbb{L}_\Theta
    &=-\ii(I\otimes H_\Theta-H_\Theta^T\otimes I) \\
    &\quad+\sum_k\gamma_k\left(L_k^*\otimes L_k-\frac12 I\otimes L_k^\dagger L_k-\frac12 (L_k^\dagger L_k)^T\otimes I\right).
    \label{eq:dense_liouvillian}
\end{align}
This representation is useful for small systems, but it stores $d^4$ complex scalars. In contrast, direct evaluation of~\eqref{eq:gksl_rates} stores $H$ and the $L_k$ matrices, requiring $O(md^2)$ storage.

\subsection{Physics-informed learning}
\label{sec:pinn_background}

A PINN approximates a trajectory or field by a neural network and penalizes the governing-equation residual~\cite{raissi2019physics,karniadakis2021physics}. For an open quantum system, a generic residual objective is
\begin{equation}
    \calJ(\phi,\Theta)=
    \lambda_{\rm data}\calJ_{\rm data}
    +\lambda_{\rm phys}\calJ_{\rm phys}
    +\lambda_0\calJ_0.
    \label{eq:pinn_general_loss}
\end{equation}
The data term compares predicted observables with measurements, the physics term penalizes the master-equation residual, and the initial-condition term anchors $\rho(0)$. Neural ODEs, DeepONets, and Fourier neural operators provide related operator-learning approaches~\cite{chen2018neural,lu2021learning,li2021fourier}. Hamiltonian, Lagrangian, and other structure-preserving networks show that hard physical constraints can improve scientific learning by restricting the hypothesis class~\cite{greydanus2019hamiltonian,cranmer2020lagrangian,hairer2006geometric}.

Quantum machine-learning work has demonstrated neural state representations, quantum-state tomography, variational steady-state ansatzes, and non-Markovian dynamics models~\cite{carleo2017solving,torlai2018neural,yoshioka2020variational,banchi2018modelling,luchnikov2020machine}. The present work is narrower: it focuses on Markovian density-matrix trajectories and asks that both the trajectory surrogate and the generator parameterization remain physically valid by construction.

\subsection{Why soft penalties are not enough}
\label{sec:soft_penalties}

Suppose a neural network outputs an arbitrary Hermitian trace-one matrix. The Bloch-sphere analogy shows the issue immediately: a qubit matrix
\begin{equation}
    \rho(r)=\frac12(I+r_x\sigma_x+r_y\sigma_y+r_z\sigma_z)
\end{equation}
is a density matrix if and only if $\norm{r}_2\le1$. A neural network with unconstrained $r(t)\in\R^3$ can leave the ball even when the loss eventually pushes it back. During optimization, such nonphysical values can corrupt residuals, produce negative probabilities, and interact badly with parameter estimation. Projection after each forward pass solves physicality only after a discontinuous or non-smooth operation. The Cholesky map instead defines the network image inside $\calD(\hilb)$.

% ============================================================
\section{Method: CPTP-Compiler-PINNs}
\label{sec:method}

\subsection{Hard-constrained density parameterization}
\label{sec:density_parameterization}

Let $B_\phi:[0,T]\to\C^{d\times d}$ be a neural network whose output is interpreted as a lower-triangular matrix. We use a nonzero diagonal to avoid a zero denominator:
\begin{equation}
    A_\phi(t)_{ij}=0\quad(i<j),\qquad
    A_\phi(t)_{ii}=a_{\min}+\softplus(b_{ii}(t)),\qquad
    A_\phi(t)_{ij}=b_{ij}(t)\quad(i>j),
    \label{eq:A_definition}
\end{equation}
where $a_{\min}>0$ is a small numerical floor. The density output is
\begin{equation}
    \rhophi(t)=\frac{A_\phi(t)A_\phi(t)^\dagger}{\Tr(A_\phi(t)A_\phi(t)^\dagger)}.
    \label{eq:cholesky_density}
\end{equation}
The lower-triangular structure is not essential for physicality, but it removes redundant degrees of freedom and connects the representation to Cholesky factorization.

The output dimension is $d^2$ real degrees for a complex lower-triangular matrix with real positive diagonal: $d$ real diagonal entries and $d(d-1)$ real off-diagonal degrees. A full-rank density matrix has $d^2-1$ real degrees after trace normalization, so the representation is only mildly redundant. Rank-deficient states are obtained as limits when some diagonal entries approach the floor and the off-diagonal factors align accordingly.

\subsection{Generator parameterization}
\label{sec:generator_parameterization}

For known jump operators with unknown rates, we write
\begin{equation}
    \gamma_k(\xi_k)=\softplus(\xi_k)+\gamma_{\min},
    \qquad \gamma_{\min}\ge0.
    \label{eq:softplus_rates}
\end{equation}
For a full Kossakowski matrix in a fixed operator basis $\{F_a\}_{a=1}^{q}$, we use
\begin{equation}
    C_\eta=B_\eta B_\eta^\dagger\succeq0,
    \label{eq:kossakowski_psd}
\end{equation}
and write
\begin{equation}
    \calL_\Theta[\rho]
    =-\ii[H_\Theta,\rho]
    +\sum_{a,b=1}^{q}(C_\eta)_{ab}
    \left(F_a\rho F_b^\dagger-\frac12\{F_b^\dagger F_a,\rho\}\right).
    \label{eq:kossakowski_generator}
\end{equation}
Hamiltonians are parameterized in a Hermitian basis $\{G_\ell\}$:
\begin{equation}
    H_\theta(t)=\sum_{\ell}\theta_\ell(t)G_\ell,
    \qquad G_\ell=G_\ell^\dagger.
\end{equation}
Together, these choices ensure that the learned generator remains in GKSL form.

\subsection{Observation model and loss}
\label{sec:loss}

Let $\{O_j\}_{j=1}^{p}$ be measured observables, let $y_{ij}$ be a measurement at time $t_i$, and let $m_{ij}\in\{0,1\}$ indicate whether the observation is present. The data loss is
\begin{equation}
    \calJ_{\rm data}(\phi)=
    \frac{1}{\sum_{i,j}m_{ij}}
    \sum_{i,j}m_{ij}\left(\Tr[O_j\rhophi(t_i)]-y_{ij}\right)^2.
    \label{eq:data_loss}
\end{equation}
The physics loss is
\begin{equation}
    \calJ_{\rm phys}(\phi,\Theta)=
    \frac{1}{N_c}\sum_{\ell=1}^{N_c}
    \norm{\dot\rhophi(\tau_\ell)-\calL_\Theta[\rhophi(\tau_\ell)]}_F^2.
    \label{eq:physics_loss}
\end{equation}
The derivative $\dot\rhophi$ may be computed by automatic differentiation. In the CPU experiments reported here we used the centered finite-difference surrogate
\begin{equation}
    \dot\rhophi(t)\approx\frac{\rhophi(t+h)-\rhophi(t-h)}{2h}
    \label{eq:finite_difference_derivative}
\end{equation}
for speed and reproducibility. The initial-condition loss is
\begin{equation}
    \calJ_0(\phi)=\norm{\rhophi(0)-\rho_0}_F^2.
    \label{eq:ic_loss}
\end{equation}
The full training problem is the minimization of~\eqref{eq:pinn_general_loss}. We used Adam in the accompanying CPU experiments~\cite{kingma2015adam}; L-BFGS is a natural second-stage optimizer~\cite{liu1989limited}.

\subsection{Open-system intermediate representation}
\label{sec:ir}

The compiler receives an IR
\begin{equation}
    \mathsf{IR}=\left(d,\ \{G_\ell\},\ \{F_a\},\ \Theta,\ \{O_j\},\ \mathsf{batch},\ \mathsf{mode}\right).
\end{equation}
The mode is either dense or structured. Dense mode constructs~\eqref{eq:dense_liouvillian}. Structured mode evaluates~\eqref{eq:gksl_rates} directly. The structured path is the default when $d$ is large, $m$ is small, or the computation is batched across many time points or parameter samples. In a GPU implementation, each collocation point can be mapped to a thread block or a batched GEMM tile, and jump-operator terms can be fused.


\subsection{Compiler dispatch policy}
\label{sec:dispatch_policy}

The compiler uses three layers of decisions. The first is \emph{representation}: whether to use dense superoperators, direct matrix products, sparse matrices, or tensor-local kernels. The second is \emph{batching}: whether to batch over time points, initial states, parameter samples, or measurement shots. The third is \emph{differentiation}: whether gradients are obtained by automatic differentiation through every residual evaluation, by adjoint-state equations, or by hybrid finite differences for selected nuisance parameters.

A practical dispatch rule is
\begin{equation}
    \mathsf{mode}=\begin{cases}
    \mathsf{dense}, & d\le d_{\rm dense}\text{ and }N_{\rm batch}\text{ is small},\\
    \mathsf{structured}, & md^2\ll d^4\text{ or }N_{\rm batch}\text{ is large},\\
    \mathsf{tensor}, & H\text{ and }L_k\text{ are sums of tensor-local terms}.
    \end{cases}
\end{equation}
The thresholds should be empirical because BLAS libraries, GPU memory, and batch sizes change the crossover. The compiler benchmark in Section~\ref{sec:experiment_compiler} illustrates this: dense mode is faster for tiny matrices, while structured mode becomes superior as $d$ increases.

\begin{algorithm}[t]
\caption{Residual-kernel dispatch from an open-system IR}
\label{alg:dispatch}
\begin{algorithmic}[1]
\REQUIRE IR with dimension $d$, jump count $m$, operator sparsity metadata, tensor-local metadata, batch size $N_b$, and target device.
\IF{tensor-local metadata is available and tensor contraction backend is enabled}
    \STATE return tensor-local residual kernel.
\ELSIF{$d^4$ storage fits device memory and $d\le d_{\rm dense}$}
    \STATE materialize dense Liouvillian and return dense matvec kernel.
\ELSIF{operators are sparse}
    \STATE return sparse structured kernel using sparse-dense products.
\ELSE
    \STATE return dense-matrix structured kernel evaluating commutator and dissipators term by term.
\ENDIF
\STATE Attach gradient rule: automatic differentiation, custom adjoint, or finite-difference fallback.
\end{algorithmic}
\end{algorithm}

This dispatch design is intentionally close to accelerator compilation: the IR preserves mathematical structure long enough for the backend to decide how to lower it. A dense Liouvillian is only one lowering, not the canonical representation of the physics.

\subsection{Algorithm}
\label{sec:algorithm}

\begin{algorithm}[t]
\caption{CPTP-Compiler-PINN training}
\label{alg:cptp_compiler_pinn}
\begin{algorithmic}[1]
\REQUIRE Observations $\{(t_i,O_j,y_{ij},m_{ij})\}$, initial state $\rho_0$, Hamiltonian basis $\{G_\ell\}$, jump basis $\{L_k\}$ or Kossakowski basis $\{F_a\}$, collocation grid $\{\tau_\ell\}$.
\ENSURE Density surrogate $\rhophi(t)$ and generator parameters $\Theta$.
\STATE Initialize neural parameters $\phi$ and unconstrained generator parameters $\xi$.
\STATE Compile the open-system IR to a residual evaluator $\mathsf{evalRHS}$ using dense or structured mode.
\FOR{training step $s=1,2,\ldots,S$}
    \STATE Form $A_\phi(t)$ by~\eqref{eq:A_definition} and $\rhophi(t)$ by~\eqref{eq:cholesky_density}.
    \STATE Map unconstrained rates to $\gamma_k=\softplus(\xi_k)+\gamma_{\min}$, or set $C=B B^\dagger$ for a Kossakowski model.
    \STATE Evaluate observable predictions $\Tr[O_j\rhophi(t_i)]$.
    \STATE Evaluate $\dot\rhophi(\tau_\ell)$ by automatic differentiation or~\eqref{eq:finite_difference_derivative}.
    \STATE Evaluate $\calL_\Theta[\rhophi(\tau_\ell)]$ with $\mathsf{evalRHS}$.
    \STATE Compute $\calJ=\lambda_{\rm data}\calJ_{\rm data}+\lambda_{\rm phys}\calJ_{\rm phys}+\lambda_0\calJ_0$.
    \STATE Update $(\phi,\xi)$ with Adam or a quasi-Newton optimizer.
\ENDFOR
\STATE Return $(\phi,\Theta)$ and the residual certificate from Theorem~\ref{thm:certificate}.
\end{algorithmic}
\end{algorithm}

% ============================================================
\section{Theory}
\label{sec:theory}

This section collects the theoretical guarantees. Proofs are included inline for readability and are expanded in Appendix~\ref{sec:appendix_proofs} where additional implementation details are needed.

\subsection{Exact physicality}
\label{sec:physicality_theory}

\begin{theorem}[Exact density-matrix physicality]
\label{thm:physicality}
Let $A:[0,T]\to\C^{d\times d}$ be any continuous matrix-valued function such that $A(t)\ne0$ for all $t$. Define
\begin{equation}
    \rho_A(t)=\frac{A(t)A(t)^\dagger}{\Tr(A(t)A(t)^\dagger)}.
\end{equation}
Then $\rho_A(t)\in\calD(\hilb)$ for every $t\in[0,T]$. If $A$ is differentiable, then $\rho_A$ is differentiable, Hermitian for every $t$, and $\Tr\dot\rho_A(t)=0$.
\end{theorem}

\begin{proof}
Hermiticity follows from $(AA^\dagger)^\dagger=AA^\dagger$. For every $v\in\C^d$, $v^\dagger AA^\dagger v=\norm{A^\dagger v}_2^2\ge0$, so $AA^\dagger\succeq0$. The denominator is positive because $A\ne0$ implies $\Tr(AA^\dagger)=\norm{A}_F^2>0$. Normalization gives unit trace. Differentiability follows from the quotient rule on a nonzero denominator. Differentiating $\Tr\rho_A(t)=1$ gives $\Tr\dot\rho_A(t)=0$.
\end{proof}

The theorem is simple, but it is operationally important: no training step can output a negative-probability state unless the floating-point implementation itself fails.

\subsection{Universal approximation on density trajectories}
\label{sec:universal_theory}

The original draft assumed strictly positive states. The boundary of $\calD(\hilb)$ matters because pure states and leakage-free subspaces are rank deficient. The next theorem handles continuous trajectories that may touch the boundary.

\begin{theorem}[Universal approximation of continuous density trajectories]
\label{thm:universal}
Let $\rho:[0,T]\to\calD(\hilb)$ be continuous. Let $\sigma\in\calD(\hilb)$ be full rank, for example $\sigma=I/d$. For every $\eps>0$ there exists a finite neural network $A_\phi:[0,T]\to\C^{d\times d}$ with nonzero output such that
\begin{equation}
    \sup_{t\in[0,T]}\norm{\frac{A_\phi(t)A_\phi(t)^\dagger}{\Tr[A_\phi(t)A_\phi(t)^\dagger]}-\rho(t)}_F<\eps.
\end{equation}
\end{theorem}

\begin{proof}
Fix $\eta\in(0,1)$ and define the full-rank path
\begin{equation}
    \rho_\eta(t)=(1-\eta)\rho(t)+\eta\sigma.
\end{equation}
Then $\rho_\eta(t)\succ0$ for every $t$, and
\begin{equation}
    \sup_t\norm{\rho_\eta(t)-\rho(t)}_F=\eta\sup_t\norm{\sigma-\rho(t)}_F.
\end{equation}
Choose $\eta$ small enough that this term is below $\eps/2$. Because $\rho_\eta(t)$ is continuous and uniformly positive definite on compact $[0,T]$, its Cholesky factor $C_\eta(t)$ can be chosen continuous with positive diagonal. The real and imaginary lower-triangular entries define a continuous map from $[0,T]$ to $\R^{d^2}$. Standard universal approximation theorems for feedforward networks imply that a finite network can approximate this map uniformly~\cite{cybenko1989approximation,hornik1991approximation}. The map $A\mapsto AA^\dagger/\Tr(AA^\dagger)$ is locally Lipschitz on any compact set bounded away from zero. Therefore a sufficiently accurate approximation of $C_\eta$ yields an approximation of $\rho_\eta$ within $\eps/2$. The triangle inequality completes the proof.
\end{proof}

\begin{remark}
The theorem is an approximation result, not a training guarantee. It states that the architecture does not exclude any continuous physical trajectory, up to arbitrary precision. Optimization remains nonconvex.
\end{remark}

\subsection{CPTP generator preservation}
\label{sec:cptp_generator_theory}

\begin{theorem}[Softplus Lindblad parameterization preserves GKSL form]
\label{thm:cptp_generator}
Let $H_\theta=H_\theta^\dagger$ and let $\gamma_k(\xi_k)=\softplus(\xi_k)+\gamma_{\min}$ with $\gamma_{\min}\ge0$. Then the generator~\eqref{eq:gksl_rates} is a GKSL generator and $e^{t\calL_\Theta}$ is CPTP for every $t\ge0$. Likewise, the Kossakowski parameterization~\eqref{eq:kossakowski_generator} with $C=BB^\dagger$ is a GKSL generator.
\end{theorem}

\begin{proof}
The softplus function is nonnegative, so all rates are nonnegative. The Hamiltonian is Hermitian by construction. Thus~\eqref{eq:gksl_rates} is exactly in GKSL form, and the GKSL theorem implies that it generates a CPTP semigroup~\cite{lindblad1976generators,gorini1976completely}. For~\eqref{eq:kossakowski_generator}, $C=BB^\dagger$ is positive semidefinite. Diagonalizing $C=U\Lambda U^\dagger$ and defining transformed jump operators $\widetilde L_r=\sum_a U_{ar}\sqrt{\Lambda_r}F_a$ rewrites the dissipator in standard GKSL form.
\end{proof}

This theorem concerns the learned generator. The Cholesky density network separately guarantees physical predicted states even before the residual is small.

\subsection{Trace-norm residual certificate}
\label{sec:certificate_theory}

A useful feature of residual learning is that a small residual should certify a small trajectory error. The next theorem gives a dimension-explicit, norm-stable statement in trace norm. It is sharper than a generic Frobenius Gronwall bound because CPTP maps contract trace distance.

\begin{theorem}[A posteriori residual certificate]
\label{thm:certificate}
Let $\rho(t)$ solve $\dot\rho=\calL[\rho]$ with $\rho(0)=\rho_0$, where $\calL$ generates a CPTP semigroup $\Phi_t=e^{t\calL}$. Let $\rhohat(t)$ be a differentiable Hermitian, trace-one path, and define the residual
\begin{equation}
    r(t)=\dot\rhohat(t)-\calL[\rhohat(t)].
\end{equation}
Then for every $t\in[0,T]$,
\begin{equation}
    \frac12\norm{\rhohat(t)-\rho(t)}_1
    \le
    \frac12\norm{\rhohat(0)-\rho_0}_1
    +\frac12\int_0^t\norm{r(s)}_1\,\dd s.
    \label{eq:certificate}
\end{equation}
\end{theorem}

\begin{proof}
Let $e(t)=\rhohat(t)-\rho(t)$. Variation of constants gives
\begin{equation}
    e(t)=\Phi_t[e(0)]+\int_0^t\Phi_{t-s}[r(s)]\,\dd s.
\end{equation}
The path $\rhohat$ is Hermitian and trace-one, and $\calL$ is trace preserving; hence $r(s)$ is Hermitian and traceless. A CPTP map contracts trace distance, and by homogeneity this implies $\norm{\Phi_t[X]}_1\le\norm{X}_1$ for Hermitian traceless $X$. Applying the triangle inequality yields~\eqref{eq:certificate}.
\end{proof}

The result can be evaluated from a trained model: compute the trace norm of the residual along a validation grid and integrate it. It is conservative, but it has a clear operational meaning.

\subsection{Local identifiability}
\label{sec:identifiability_theory}

No architecture can identify parameters that are invisible to the selected measurements. The following theorem states the standard local condition in a form adapted to Lindblad residual learning.

\begin{theorem}[Local parameter identifiability through the sensitivity Gramian]
\label{thm:identifiability}
Let $\Theta\in\R^q$ parameterize a smooth GKSL generator, and let
\begin{equation}
    f_{ij}(\Theta)=\Tr[O_j\rho_\Theta(t_i)]
\end{equation}
be the noiseless observation map for fixed initial state and measurement set. Let $J(\Theta_*)\in\R^{M\times q}$ be the Jacobian of all observed entries at the true parameter $\Theta_*$. If $J(\Theta_*)$ has full column rank with smallest singular value $s_{\min}>0$, then there exist neighborhoods $U$ of $\Theta_*$ and $V$ of $f(\Theta_*)$ such that, for any estimate $\Thetahat\in U$ with observation residual $\norm{f(\Thetahat)-f(\Theta_*)}_2\le\delta$, one has
\begin{equation}
    \norm{\Thetahat-\Theta_*}_2\le \frac{2\delta}{s_{\min}}
\end{equation}
provided $U$ is small enough that the Jacobian varies by at most $s_{\min}/2$ in operator norm.
\end{theorem}

\begin{proof}
By the fundamental theorem of calculus,
\begin{equation}
    f(\Thetahat)-f(\Theta_*)=\left(\int_0^1 J(\Theta_*+\alpha(\Thetahat-\Theta_*))\,\dd\alpha\right)(\Thetahat-\Theta_*).
\end{equation}
The averaged Jacobian has smallest singular value at least $s_{\min}/2$ by the perturbation assumption. Therefore
\begin{equation}
    \delta\ge\norm{f(\Thetahat)-f(\Theta_*)}_2\ge \frac{s_{\min}}{2}\norm{\Thetahat-\Theta_*}_2.
\end{equation}
\end{proof}

The theorem also explains why some inverse problems remain ill-conditioned even with hard physical constraints. If the measurement set does not excite a rate or Hamiltonian term, the sensitivity Gramian is singular.

\subsection{Compiler complexity}
\label{sec:compiler_theory}

\begin{theorem}[Dense versus structured residual complexity]
\label{thm:compiler}
Consider a $d$-dimensional system with $m$ dense Lindblad operators. Dense Liouvillian storage requires $O(d^4)$ memory, and explicit construction costs $O(md^4)$ arithmetic. A dense residual evaluation by matrix-vector multiplication costs $O(d^4)$ per state. Structured evaluation of~\eqref{eq:gksl_rates} costs $O(md^3)$ arithmetic and $O(md^2)$ storage. For $N$ collocation states, dense evaluation after construction costs $O(md^4+Nd^4)$, while structured evaluation costs $O(Nmd^3)$.
\end{theorem}

\begin{proof}
The dense superoperator has dimension $d^2\times d^2$, hence $d^4$ entries. Each jump contributes Kronecker products of $d\times d$ matrices, giving $O(d^4)$ entries per jump. Multiplication of a dense $d^2\times d^2$ matrix by a vector costs $O(d^4)$. Structured evaluation uses matrix products $H\rho$, $\rho H$, $L_k\rho$, $(L_k\rho)L_k^\dagger$, $(L_k^\dagger L_k)\rho$, and $\rho(L_k^\dagger L_k)$; for dense matrices each is $O(d^3)$, repeated $m$ times. The storage is one Hamiltonian plus $m$ jump operators, each $d^2$ entries.
\end{proof}

\begin{remark}
If the $L_k$ are sparse or tensor-local, the structured cost can be lower than $O(md^3)$. If the dense superoperator is reused over extremely large batches, dense matvecs can be competitive for small $d$. The compiler should therefore benchmark and dispatch by regime rather than enforce a single mode.
\end{remark}


\subsection{Sensitivity equations for experiment design}
\label{sec:sensitivity_equations}

The identifiability theorem can be made computational by differentiating the master equation. Let $S_a(t)=\partial \rho_\Theta(t)/\partial\Theta_a$. If the initial state does not depend on $\Theta$, then
\begin{equation}
    \dot S_a(t)=\calL_\Theta[S_a(t)]+\left(\partial_{\Theta_a}\calL_\Theta\right)[\rho_\Theta(t)],
    \qquad S_a(0)=0.
    \label{eq:sensitivity_ode}
\end{equation}
The observation Jacobian in Theorem~\ref{thm:identifiability} has entries
\begin{equation}
    J_{(i,j),a}=\Tr[O_j S_a(t_i)].
\end{equation}
Thus the same compiler used for residual evaluation can evaluate sensitivity equations. This suggests an active experimental design loop: choose preparation states, observables, and time points that maximize the smallest singular value of the sensitivity matrix or minimize a posterior covariance. Such a loop is outside the scope of the reported simulations, but it is a direct extension of the theory.

\subsection{Low-rank and purification variants}
\label{sec:low_rank_variants}

For larger Hilbert spaces, a full $d\times d$ Cholesky factor may be too expensive. A rank-$r$ variant uses $A_\phi(t)\in\C^{d\times r}$ and
\begin{equation}
    \rho_\phi^{(r)}(t)=\frac{A_\phi(t)A_\phi(t)^\dagger+\alpha I/d}{\Tr[A_\phi(t)A_\phi(t)^\dagger]+\alpha},
    \qquad \alpha\ge0.
    \label{eq:low_rank_density}
\end{equation}
When $\alpha>0$, the state is full rank; when $\alpha=0$, its rank is at most $r$. This variant connects to purification-based simulation: $A$ can be viewed as an unnormalized purification amplitude with an ancilla of dimension $r$. The same physicality theorem applies. Universal approximation of arbitrary full-rank states requires $r=d$ or an isotropic floor; approximation of low-rank trajectories can use smaller $r$.

\begin{corollary}[Low-rank physicality]
\label{cor:low_rank}
For any $A\in\C^{d\times r}$ and $\alpha\ge0$ not both zero, the state in~\eqref{eq:low_rank_density} is Hermitian, positive semidefinite, and unit trace.
\end{corollary}

\begin{proof}
Both $AA^\dagger$ and $\alpha I/d$ are positive semidefinite. Their sum is nonzero by assumption, and trace normalization gives a density matrix.
\end{proof}

% ============================================================
\section{Numerical experiments}
\label{sec:experiments}

\subsection{Reproducibility and scope}
\label{sec:experimental_scope}

All numerical values in this section were generated by the scripts in the accompanying folder. The simulations were run on CPU using PyTorch and SciPy~\cite{paszke2019pytorch,virtanen2020scipy}. The qubit PINN experiments used Adam~\cite{kingma2015adam}. The goal is to validate structural claims, not to claim production-scale performance. In particular, the qubit training uses a finite-difference residual rather than second-order automatic differentiation because the purpose of the artifact is reproducibility in a lightweight environment. Larger studies should use batched automatic differentiation, mixed precision where safe, and GPU-resident residual kernels.

\subsection{Experiment 1: qubit inverse problem}
\label{sec:experiment_qubit}

We considered a driven qubit with amplitude damping and pure dephasing:
\begin{align}
    H &= \frac{\omega}{2}\sigma_z+\frac{\Omega}{2}\sigma_x,\\
    \calL[\rho] &= -\ii[H,\rho]
    +\gamma_1\left(\sigma_-\rho\sigma_+-\frac12\{\sigma_+\sigma_-,\rho\}\right)
    +\frac{\gamma_\phi}{2}\left(\sigma_z\rho\sigma_z-\rho\right).
\end{align}
The ground-truth parameters were
\begin{equation}
    \omega=\pi,
    \qquad \Omega=0.3\pi,
    \qquad \gamma_1=0.35,
    \qquad \gamma_\phi=0.18.
\end{equation}
The initial state was $\rho_0=\ket{+}\bra{+}$, and observations were noisy samples of $\langle\sigma_x\rangle$, $\langle\sigma_y\rangle$, and $\langle\sigma_z\rangle$ on a grid of 41 times. Gaussian noise with standard deviation $0.01$ was added. We compared the CPTP-Cholesky network with a Hermitian trace-one baseline of the form $\rho(t)=\frac12(I+x(t)\sigma_x+y(t)\sigma_y+z(t)\sigma_z)$, which does not enforce positivity.

\begin{figure}[H]
    \centering
    \includegraphics[width=0.98\textwidth]{fig_qubit_observables.pdf}
    \caption{Qubit inverse problem with 50\% of the time grid observed. The CPTP-Cholesky model and Hermitian trace-one baseline are trained from the same noisy observations. The hard-constrained model stays inside the density-matrix set by construction.}
    \label{fig:qubit_observables}
\end{figure}

\input{tables/exp1_parameter_values.tex}

Table~\ref{tab:parameter_values} gives one representative full-observation run. The parameter estimates are not exact; the training budget was deliberately small. The relevant point is that the constrained model obtains a lower trajectory error while preserving positivity. Table~\ref{tab:qubit_inverse} summarizes two seeds at three observation fractions. In these runs, the CPTP-Cholesky network had lower mean trace distance than the baseline at all observation fractions. The baseline often remained close to the positive cone, but its minimum eigenvalues approached or crossed zero because positivity was not guaranteed.

\input{tables/exp1_qubit_inverse.tex}

\begin{figure}[H]
    \centering
    \includegraphics[width=0.58\textwidth]{fig_sparse_ablation.pdf}
    \caption{Sparse-observation ablation for the qubit inverse problem. Error bars are sample standard deviations over two seeds. This plot is an authentic small-scale simulation, not a broad benchmark claim.}
    \label{fig:sparse_ablation}
\end{figure}

\begin{figure}[H]
    \centering
    \includegraphics[width=0.62\textwidth]{fig_physicality_box.pdf}
    \caption{Minimum eigenvalue over predicted qubit trajectories. The CPTP-Cholesky model is positive by construction; the Hermitian trace-one baseline has no architectural positivity guarantee.}
    \label{fig:physicality_box}
\end{figure}

\subsection{Experiment 2: residual certificate}
\label{sec:experiment_certificate}

The residual theorem was tested by taking the exact qubit trajectory $\rho(t)$ and forming physical perturbations
\begin{equation}
    \rho_\eps(t)=(1-a_\eps(t))\rho(t)+a_\eps(t)I/2,
    \qquad a_\eps(t)=\eps\sin^2(\pi t/T).
\end{equation}
These perturbed trajectories remain valid density matrices. We evaluated the residual $r_\eps(t)=\dot\rho_\eps(t)-\calL[\rho_\eps(t)]$ and compared the actual trace distance with the certificate in Theorem~\ref{thm:certificate}. The certificate is conservative by a nearly constant factor in this test, as shown in Table~\ref{tab:residual_certificate} and Figure~\ref{fig:residual_certificate}.

\input{tables/exp2_residual_certificate.tex}

\begin{figure}[H]
    \centering
    \includegraphics[width=0.60\textwidth]{fig_residual_certificate.pdf}
    \caption{Trace-norm residual certificate versus actual error for physical perturbations of the exact qubit trajectory. The certificate scales linearly with perturbation amplitude and upper bounds the observed trace-distance error.}
    \label{fig:residual_certificate}
\end{figure}

\subsection{Experiment 3: qutrit leakage identification}
\label{sec:experiment_qutrit}

A transmon-like leakage toy model was used to validate positive-rate identification outside the qubit case. The basis states are $\ket{0},\ket{1},\ket{2}$, with cascade decay $\ket{2}\to\ket{1}\to\ket{0}$:
\begin{equation}
    L_{10}=\ket{0}\bra{1},\qquad L_{21}=\ket{1}\bra{2}.
\end{equation}
The true rates were $\gamma_{10}=0.24$ and $\gamma_{21}=0.075$. Population measurements were generated from the exact cascade dynamics, perturbed by small noise, clipped to $[0,1]$, and renormalized. A softplus-positive generator fit recovered both rates accurately in this controlled setting.

\input{tables/exp3_qutrit_leakage.tex}

\begin{figure}[H]
    \centering
    \includegraphics[width=0.66\textwidth]{fig_qutrit_leakage.pdf}
    \caption{Qutrit leakage identification from population observations. Dots are noisy population samples; curves show exact and fitted population dynamics.}
    \label{fig:qutrit_leakage}
\end{figure}

\subsection{Experiment 4: compiler scaling}
\label{sec:experiment_compiler}

The compiler experiment compares dense superoperator evaluation with structured matrix-free residual evaluation on random dense Lindblad systems. For small dimensions, dense matrix-vector products can be faster because the superoperator is tiny and BLAS overhead is low. As $d$ increases, the storage and evaluation costs of dense mode grow rapidly. In the CPU benchmark, structured evaluation became faster by $d=8$ and used far less memory for all nontrivial dimensions. The equality of the two residuals was verified to floating-point precision.

\input{tables/exp4_compiler_scaling.tex}

\begin{figure}[H]
    \centering
    \includegraphics[width=0.94\textwidth]{fig_compiler_scaling.pdf}
    \caption{Compiler scaling benchmark. Left: dense superoperator storage versus structured IR storage. Right: residual evaluation time. The small-$d$ regime favors dense mode; the larger-$d$ regime favors structured mode.}
    \label{fig:compiler_scaling}
\end{figure}

% ============================================================
\section{Discussion}
\label{sec:discussion}

\subsection{What is genuinely gained by the hard constraint?}

The Cholesky map changes the feasible set. A soft penalty for positivity competes with data fit, residual fit, and optimizer noise. A hard parameterization removes that competition. In the qubit experiment, the baseline sometimes came close to physicality, but it was never guaranteed. In larger systems, where the positive cone is a small subset of Hermitian trace-one matrices, the gap between hard constraints and soft penalties is expected to become more important.

The tradeoff is optimization geometry. The Cholesky map is nonlinear and redundant. Near rank-deficient states, small changes in the factor can produce small eigenvalues and ill-conditioning. This is not a reason to abandon the representation; it is a reason to use careful initialization, diagonal floors, learning-rate schedules, and possibly low-rank-plus-isotropic variants.

\subsection{Why the compiler matters}

The compiler is not merely a software convenience. It determines whether the residual can be evaluated without forming an object whose size scales as $d^4$. In a multi-qubit Hilbert space, $d=2^n$, so dense Liouvillian storage scales as $16^n$. Structured kernels still face exponential Hilbert-space growth, but they avoid an unnecessary square of the density-vector dimension. This is the difference between a prototype that only runs on toy models and a path that can exploit locality, batching, sparsity, tensor products, and GPU memory hierarchy.

The most direct accelerator path is to lower each Hamiltonian and dissipator term to batched matrix multiplications, fuse anticommutator terms, and keep density matrices resident on GPU. CUDA-Q and related quantum programming frameworks provide a natural place for this IR to interact with circuit-level or pulse-level descriptions~\cite{cudaq2023,qiskit2024}. A research scientist or compiler engineer would likely focus next on kernel fusion, auto-tuning, sparse/tensor-local jump operators, mixed precision with physicality checks, and adjoint differentiation.

\subsection{Relationship to quantum tomography and neural quantum states}

Quantum process tomography reconstructs channels without assuming a low-parameter GKSL model, but its resource cost is high~\cite{chuang1997prescription,poyatos1997complete}. Neural quantum states and neural tomography represent states or measurement distributions with flexible models~\cite{carleo2017solving,torlai2018neural}. The present approach occupies a different point in the design space: it assumes a Markovian master-equation form and uses the residual to couple sparse observations to a physical trajectory. When the Markovian assumption is wrong, process tensors or non-Markovian recurrent models are more appropriate~\cite{wolf2008assessing,banchi2018modelling,luchnikov2020machine}.

\subsection{Limitations}
\label{sec:limitations}

The current manuscript has several limitations that should be transparent before submission.
\begin{enumerate}[leftmargin=1.5em]
    \item \textbf{No hardware data.} The experiments are synthetic CPU simulations. They validate implementation and theory but do not demonstrate performance on a quantum device.
    \item \textbf{Small number of seeds.} The qubit table uses two seeds, enough for a reproducible artifact but not enough for a definitive benchmark.
    \item \textbf{Finite-difference residual in reported PINN runs.} The mathematical algorithm supports automatic differentiation. The reported CPU experiments used finite differences to avoid expensive higher-order autodiff in this environment.
    \item \textbf{Markovian assumption.} The framework assumes GKSL dynamics. Non-Markovian systems require extended state spaces, memory kernels, or process-tensor methods.
    \item \textbf{Identifiability is measurement-dependent.} Theorems cannot overcome uninformative observables. If the sensitivity Gramian is rank deficient, no optimizer can recover all parameters.
    \item \textbf{No state-of-the-art claim.} The results are not compared against all modern quantum system-identification methods. They should be read as a principled proof of concept.
\end{enumerate}

% ============================================================
\section{Conclusion}
\label{sec:conclusion}

CPTP-Compiler-PINNs combine physical state parameterization, GKSL-preserving generator learning, and matrix-free residual compilation. The main conceptual point is that physicality should be an architectural invariant rather than an after-the-fact penalty. The mathematical results show that this invariant does not destroy universal approximation, that the learned generator remains CPTP, that residuals provide an a posteriori trace-norm certificate, and that structured compilation avoids dense Liouvillian storage. The numerical experiments are deliberately modest but authentic: every figure and table is generated by reproducible scripts.

Future work should extend the experiments to larger multi-qubit systems, hardware-calibration datasets, non-Markovian embeddings, and GPU-native kernels. The most promising engineering direction is a compiler path that takes a high-level quantum noise model and emits fused, batched, differentiable kernels for accelerator execution. That direction is naturally aligned with the needs of quantum hardware characterization and high-performance scientific machine learning.

\section*{Data and code availability}
The numerical data, tables, figures, and scripts used for this manuscript are included with the submitted artifact folder. The scripts generate all reported numbers from fixed random seeds and do not contain hard-coded experimental outcomes. For journal submission, the folder should be placed in a public repository or archived with a DOI.

\section*{Acknowledgments}
Replace this paragraph with genuine acknowledgments before submission. No external funding or hardware access is claimed in this draft.

% ============================================================
\appendix

\section{Additional proofs and derivations}
\label{sec:appendix_proofs}

\subsection{Derivative of the Cholesky density map}

Let $M(t)=A(t)A(t)^\dagger$ and $Z(t)=\Tr M(t)$. Then
\begin{equation}
    \rho(t)=\frac{M(t)}{Z(t)},
    \qquad
    \dot\rho(t)=\frac{\dot M(t)}{Z(t)}-\frac{M(t)\dot Z(t)}{Z(t)^2},
    \label{eq:cholesky_derivative}
\end{equation}
where
\begin{equation}
    \dot M(t)=\dot A(t)A(t)^\dagger+A(t)\dot A(t)^\dagger,
    \qquad
    \dot Z(t)=\Tr\dot M(t).
\end{equation}
This expression makes it clear that $\dot\rho$ is Hermitian and traceless. In automatic differentiation frameworks, the same derivative is obtained by differentiating through the factor map and normalization. In finite precision, symmetrizing $\dot\rho\leftarrow(\dot\rho+\dot\rho^\dagger)/2$ can be used as a numerical safeguard.

\subsection{Trace-norm contraction for Hermitian traceless residuals}

The proof of Theorem~\ref{thm:certificate} uses a standard consequence of CPTP contractivity. If $X=X^\dagger$ and $\Tr X=0$, write its Jordan decomposition as $X=X_+-X_-$, where $X_+,X_-\succeq0$ and $X_+X_-=0$. Since $\Tr X=0$, $\Tr X_+=\Tr X_-=a$. If $a=0$, the claim is trivial. Otherwise $\rho_+=X_+/a$ and $\rho_-=X_-/a$ are density matrices and
\begin{equation}
    X=a(\rho_+-\rho_-),
    \qquad \norm{X}_1=2a.
\end{equation}
For any CPTP map $\Phi$, trace distance contracts:
\begin{equation}
    \frac12\norm{\Phi(\rho_+)-\Phi(\rho_-)}_1
    \le
    \frac12\norm{\rho_+-\rho_-}_1.
\end{equation}
Multiplying by $a$ gives $\norm{\Phi(X)}_1\le\norm{X}_1$.

\subsection{A Frobenius alternative}

If one prefers Frobenius norm, let $\Lambda$ be an operator norm bound for $\calL$ acting on matrices with Frobenius norm. If
\begin{equation}
    \norm{\dot\rhohat(t)-\calL[\rhohat(t)]}_F\le\delta
\end{equation}
and $\norm{\calL[X]}_F\le\Lambda\norm{X}_F$, then Gronwall's inequality gives
\begin{equation}
    \norm{\rhohat(t)-\rho(t)}_F
    \le
    e^{\Lambda t}\norm{\rhohat(0)-\rho(0)}_F
    +\frac{e^{\Lambda t}-1}{\Lambda}\delta.
\end{equation}
The trace-norm certificate is usually more meaningful for quantum states, but the Frobenius form can be easier to estimate during training.

\section{Experimental implementation details}
\label{sec:appendix_experiments}

\subsection{Qubit finite-difference PINN}

The Cholesky network used two hidden layers of width 32 with hyperbolic tangent activations. The output contained four real values: two diagonal logits and the real and imaginary parts of the lower off-diagonal entry. Diagonal entries were transformed by softplus and shifted by $10^{-4}$. The baseline used the same hidden width and depth but output a three-component Bloch vector with no constraint $\norm{r}\le1$.

The loss weights were $\lambda_{\rm data}=10$, $\lambda_{\rm phys}=1$, and $\lambda_0=100$. The optimizer was Adam with learning rate $10^{-3}$ for 1000 steps. The finite-difference step in normalized time was $10^{-3}$, and the collocation grid contained 48 points in $(0.02,0.98)$. These hyperparameters were selected for speed rather than optimality.

\subsection{Residual-certificate experiment}

The exact qubit trajectory was generated by dense matrix exponential of the true Liouvillian. Perturbations mixed the exact state with $I/2$ using $a_\eps(t)=\eps\sin^2(\pi t/T)$. The derivative was evaluated by a second-order finite-difference gradient on a 601-point grid. The certificate was computed by trapezoidal integration of $\norm{r(t)}_1/2$.

\subsection{Qutrit leakage experiment}

The qutrit population dynamics reduce to the classical cascade
\begin{equation}
    p_2(t)=e^{-\gamma_{21}t},
\end{equation}
\begin{equation}
    p_1(t)=\frac{\gamma_{21}}{\gamma_{10}-\gamma_{21}}
    \left(e^{-\gamma_{21}t}-e^{-\gamma_{10}t}\right),
    \qquad
    p_0(t)=1-p_1(t)-p_2(t),
\end{equation}
with the equal-rate limit $p_1(t)=\gamma_{21}t e^{-\gamma_{21}t}$. The reported fit used nonlinear least squares with either softplus-positive or unconstrained rates.

\subsection{Compiler benchmark}

Dense mode used the Kronecker representation in~\eqref{eq:dense_liouvillian}. Structured mode directly evaluated commutator and dissipator terms. The relative error column in Table~\ref{tab:compiler_scaling} is
\begin{equation}
    \frac{\norm{\calL_{\rm dense}[\rho]-\calL_{\rm struct}[\rho]}_F}{\max(\norm{\calL_{\rm dense}[\rho]}_F,10^{-15})}.
\end{equation}
The dense and structured evaluations agree to floating-point precision.

\section{Suggested repository structure}
\label{sec:repo_structure}

A clean repository for a public release should have the following structure:
\begin{verbatim}
ctpt-compiler-pinns/
  README.md
  environment.yml
  pyproject.toml
  src/cptp_pinns/
    density.py
    lindblad.py
    compiler.py
    losses.py
    train.py
    metrics.py
  experiments/
    exp1_qubit_inverse.py
    exp2_residual_certificate.py
    exp3_qutrit_leakage.py
    exp4_compiler_scaling.py
  figures/
  tables/
  paper/
    main.tex
    refs.bib
\end{verbatim}
Unit tests should check Hermiticity, trace, minimum eigenvalue, dense-versus-structured equivalence, gradient finiteness, reproducibility of seeds, and table regeneration. A continuous-integration job should run a small smoke test on CPU and optionally a larger benchmark on GPU.


\section{Accelerator-aware implementation roadmap}
\label{sec:accelerator_roadmap}

A production implementation should treat the mathematical model and the kernel schedule as separate layers. The model layer stores Hamiltonian terms, jump operators, rate transforms, observables, and constraints. The kernel layer chooses dense, structured, sparse, or tensor-local execution. The following roadmap is designed for GPU-oriented scientific machine learning.

\subsection{Kernel fusion}

The structured residual contains repeated products involving the same $\rho$ and operators. A naive implementation launches separate kernels for each matrix multiplication and addition. A fused implementation should combine the anticommutator terms, reuse $C_k=L_k^\dagger L_k$, and keep intermediate matrices in registers or shared memory when $d$ is small enough. For batches of trajectories, the natural primitive is strided batched GEMM. For tensor-local terms, the primitive is a sequence of small contractions rather than a full dense matrix multiply.

\subsection{Differentiation strategy}

Automatic differentiation through every residual is convenient but may store too many intermediates. Three alternatives are important: checkpointing, custom vector-Jacobian products for the Lindblad residual, and adjoint equations for sensitivities. The sensitivity equation~\eqref{eq:sensitivity_ode} is especially useful when the number of generator parameters is moderate and the number of network parameters is large. A hybrid system can use automatic differentiation for the neural density map and custom adjoints for the compiled residual.

\subsection{Precision policy}

Density matrices are sensitive to Hermiticity and trace errors. Mixed precision should therefore be used with explicit safeguards: symmetrize matrices after residual evaluation, compute traces in higher precision, monitor minimum eigenvalues on validation grids, and compare dense and structured residuals on small systems. The Cholesky parameterization protects positivity of the model output, but the residual and optimizer can still suffer from roundoff if matrix products are too low precision.

\subsection{Connection to quantum software stacks}

Quantum software frameworks such as Qiskit and CUDA-Q represent circuits, observables, and hybrid quantum-classical workflows~\cite{qiskit2024,cudaq2023}. The IR in this paper is compatible with that ecosystem: circuit- or pulse-level noise hypotheses can be compiled into Hamiltonian and jump-operator terms, while learned parameters can be exported back to simulators. For large-scale quantum advantage studies, resource-estimation and block-encoding tools may also interact with learned open-system models~\cite{beverland2022assessing,rall2023quantum}.


\section{Extended experimental protocol for a stronger submission}
\label{sec:extended_protocol}

The reported experiments are intentionally lightweight. A stronger journal submission should add a statistically powered benchmark suite. This appendix specifies such a suite in a way that preserves the central rule of this manuscript: results must be generated by executable code and must never be invented in the text.

\subsection{Protocol A: multi-seed qubit inverse benchmark}

The qubit experiment in Section~\ref{sec:experiment_qubit} should be repeated over at least 20 random seeds, five noise levels, and five observation fractions. Each seed should generate a new observation mask and measurement-noise realization. The report should include the median, interquartile range, and worst-case trace distance, plus the fraction of runs in which each method violates positivity. If the CPTP-Cholesky model is claimed to outperform a baseline, the claim should be supported by paired statistical tests on the same seeds. The baseline set should include at least: an unconstrained Hermitian trace-one PINN, a penalty-based positivity PINN, a direct ODE parameter fit, and a projection-after-forward-pass method. All models should use matched parameter budgets where possible.

\subsection{Protocol B: two-qubit dissipative calibration}

A two-qubit experiment should use a Hamiltonian with exchange coupling and local detunings,
\begin{equation}
    H(t)=J(t)(\sigma_x\otimes\sigma_x+\sigma_y\otimes\sigma_y)
    +\frac{\Delta_1}{2}\sigma_z\otimes I
    +\frac{\Delta_2}{2}I\otimes\sigma_z,
\end{equation}
with local amplitude damping and dephasing on each qubit. The training data should include multiple initial states, not only $\ket{00}$. The evaluation should report state fidelity, trace distance, recovered rates, and whether the learned generator predicts held-out initial states. A meaningful generalization test is to train on computational-basis initial states and evaluate on superposition and Bell states. The structured compiler should be compared against dense mode, and the memory footprint should be reported.

\subsection{Protocol C: qutrit leakage under partial observability}

The qutrit leakage experiment should be extended to partial observability by hiding the $\ket{2}$ population and asking whether leakage can still be inferred from distortions in the computational subspace. This is a harder and more realistic task for transmon-like devices. The sensitivity Gramian in Theorem~\ref{thm:identifiability} should be computed for each measurement design. If the Gramian is ill-conditioned, the paper should state that the parameter is weakly identifiable rather than reporting an overconfident estimate.

\subsection{Protocol D: accelerator benchmark}

A GPU benchmark should separate three timings: data transfer, residual evaluation, and backpropagation. It should report hardware model, driver version, precision, batch size, Hilbert dimension, number of jumps, sparsity, and whether kernels are fused. The dense baseline should be allowed to precompute the Liouvillian when the generator is fixed; the structured method should be allowed to precompute $L_k^\dagger L_k$. A fair benchmark should include both fixed-generator trajectory fitting and learnable-generator training, because dense precomputation is less useful when the rates or Hamiltonian change every optimization step.

\subsection{Protocol E: hardware-facing validation}

For hardware data, the paper should specify the calibration experiment, state preparations, measurement operators, readout-error mitigation, shot counts, and drift handling. The model should be trained on one time window and validated on a later time window to test temporal robustness. If the hardware data are proprietary, the paper should at minimum provide a public synthetic benchmark with the same dimensions and noise statistics. Claims about device characterization should be limited to the evidence actually reported.

\subsection{Reporting standard}

Every figure should be generated by a script that writes the underlying numerical values to disk. Every table should be generated from those values, not typed by hand. The repository should include a command such as
\begin{verbatim}
python -m experiments.run_all --output artifacts/
\end{verbatim}
that regenerates figures, tables, and a machine-readable results file. The manuscript should cite the exact commit hash used for submission. This standard is especially important for papers that are meant to demonstrate research maturity to an industrial hiring manager: reproducibility and restraint are more impressive than unsupported superlatives.


\section{Generator gauge, failure modes, and interpretation}
\label{sec:gauge_failure_modes}

\subsection{Gauge freedom in Lindblad operators}

A Lindblad representation is not unique. Unitary rotations of the jump-operator list can leave the dissipator invariant, and Hamiltonian terms can shift under more general affine transformations of the operators. Consequently, a paper should be careful when claiming recovery of individual jump operators. If the operator basis is fixed in advance, rate recovery is meaningful within that basis. If the operators are learned freely, the identifiable object is the generator or the channel family, not a unique list of physical mechanisms. This is one reason the experiments in this draft use fixed jump operators and learn only a small number of rates.

For a fixed positive Kossakowski matrix $C$, diagonalization gives many equivalent jump bases. If $C=U\Lambda U^\dagger$, then one valid jump set is
\begin{equation}
    \widetilde L_r=\sum_a U_{ar}\sqrt{\Lambda_r}F_a.
\end{equation}
Replacing $U$ by $UV$ inside a degenerate eigenspace gives a different jump list and the same generator. The compiler should therefore store both the basis-level Kossakowski representation and any diagonalized jump representation used for execution. The former is better for parameter interpretation; the latter may be better for kernels.

\subsection{Distinguishing state physicality from model correctness}

The Cholesky network guarantees that $\rho_\phi(t)$ is a valid density matrix. It does not guarantee that the trajectory is generated by the learned Lindbladian. That is the role of the residual. A model can be physical at every time and still have a large residual; such a model is a valid curve in state space but a poor solution of the proposed dynamics. Conversely, a generator can be CPTP and still fit the wrong physical mechanism if the observations are insufficient. The framework therefore has three separate diagnostics: state physicality, residual size, and prediction error on held-out observations.

This distinction is important for communication. A valid claim is: ``the architecture never outputs a nonphysical density matrix.'' A stronger claim such as ``the method discovers the true dissipative mechanism'' requires identifiability analysis, held-out validation, and comparison with alternative mechanisms. The sensitivity theorem provides one route to that analysis.

\subsection{Common failure modes}

The following failure modes are expected in practice.
\begin{enumerate}[leftmargin=1.5em]
    \item \textbf{Residual under-resolution.} If collocation points miss fast oscillations or pulses, the learned trajectory can fit data while violating dynamics between collocation points. Adaptive collocation or causal training can reduce this risk~\cite{wang2022respecting}.
    \item \textbf{Rate sloppiness.} Different combinations of Hamiltonian and dissipative parameters may produce nearly indistinguishable observables. The sensitivity Gramian will have small singular values.
    \item \textbf{Boundary ill-conditioning.} Nearly pure states can make Cholesky factors ill-conditioned. A small isotropic floor or low-rank-plus-floor representation can stabilize training.
    \item \textbf{Loss imbalance.} If the data term dominates, the residual is ignored. If the residual dominates, noisy data may be underfit. Dynamic loss weighting or nondimensionalization is often necessary.
    \item \textbf{Dense fallback misuse.} Dense mode is convenient but can silently become the bottleneck. The compiler should log memory estimates before materializing a Liouvillian.
\end{enumerate}

\subsection{Interpreting the current numerical results}

The current tables support limited conclusions only. They show that the scripts run, the hard constraint preserves positivity, the residual certificate behaves as the theorem predicts, the qutrit rate fit works in an identifiable toy model, and the structured compiler avoids dense storage. They do not prove that the method is optimal, that it scales to many qubits without tensor structure, or that it will outperform specialized tomography on every task. This restraint is intentional. A manuscript sent to a research manager should be ambitious in design but conservative in evidence; unsupported claims are easy to spot and hard to repair.

\section{Practical checklist before submission}
\label{sec:submission_checklist}

Before submitting this manuscript to a journal or sending it to a hiring manager, the following steps are necessary:
\begin{enumerate}[leftmargin=1.5em]
    \item Replace the placeholder author, affiliation, acknowledgments, and repository URL with real information.
    \item Run the experiments independently from a clean environment and verify that all numbers are regenerated.
    \item Add more seeds and stronger baselines if making any comparative-performance claim.
    \item If claiming GPU performance, run the GPU implementation and report hardware, precision, batch size, kernel choices, and profiling methodology.
    \item If claiming hardware relevance, include real device data or clearly state that the paper is simulation-only.
    \item Use the current IOP template requested by the journal at submission time, or confirm that standard LaTeX is accepted.
\end{enumerate}

% ============================================================
\bibliographystyle{unsrt}
\bibliography{refs}

\end{document}
